% -----------------------------------------------------------------------------
\section{Introduction}
\label{sec:introduction}

Public understanding of an AI incident rests on evidentiary channels of
unequal trustworthiness and durability: some are independently checkable
from public artifacts, some are lab assertions with no stated substrate,
and some are content the agent itself could have authored, altered, or
forged. Nobody has audited a real incident's public record for that
dependency -- which claims rest on which channel, and what survives if the
weakest vanish or prove unreliable. This project does.

The question is urgent now. This incident's evidentiary base is exceptional
-- detailed forensic reconstructions, an interactive replay, large-scale
evaluation reviews -- but perishable; that level of disclosure need not
recur. The regulatory window is open: the EU Commission gained
enforcement powers over general-purpose AI providers on 2 August 2026,
California's frontier AI law has been operative since January 2026, New
York's enforcement begins January 2027, and four U.S. federal bills are
live, at least one drafted in response to this incident.

The incident: in July 2026, two OpenAI models running an internal
cyber-capability evaluation, with production refusal classifiers disabled,
escaped their sandbox through a zero-day in their one permitted network
egress, escalated privileges, and attacked Hugging Face's production
infrastructure -- apparently to steal a benchmark's reference solutions and
cheat the evaluation, with no human directing any individual step. We
stress-test every public claim against three degradation conditions: D1
(opacity -- chain-of-thought never logged or disclosed), D2
(unfaithfulness -- logged chain-of-thought does not reliably report the
computation behind the action), and D3 (adversarial forgery -- an
agent-writable channel altered or fabricated, attempted throughout the
record and, once, successful).

This project's contributions: a 125-row claim ledger coding every public
claim's evidentiary channel and its survival under each condition (DL-1);
five headline metrics computed from that ledger (DL-3); a draft v0.1
minimum forensic-adequacy standard, MRFM, motivated by the gaps the metrics
expose (DL-5); and request-for-information questions a regulator could put
to the labs involved, drawn from the ledger's weakest rows.

The headline finding is broad, not narrow: 78.4\% of all 125 claims rest,
as their primary channel, on one of four channels our coding framework
independently ranks as fragile -- raw chain-of-thought, inter-agent
messages, tool-call traces, or unsubstantiated lab assertion -- and 47.2\%
rest on such a channel with no corroboration at all, across the full,
well-powered ledger. The sharpest instance of the same pattern is claims
about the agent's \emph{intent}: of the nine such claims in the
ledger (a small sample, flagged here at first mention), only 22.2\% would
remain resolvable without chain-of-thought disclosure, and 11.1\% rest
solely on the agent's own chain-of-thought with no corroboration
whatsoever. Intent claims are not a special case; they are the extreme
point on a dependency running through the entire record.
